\chapter{Разработка и реализация системы}

\section{Реализация сервиса аутентификации}

\subsection{Архитектура сервиса аутентификации}

Сервис аутентификации реализован на языке программирования Rust с использованием веб-фреймворка Axum. Сервис отвечает за регистрацию и авторизацию пользователей, управление профилем, а также выпуск и валидацию JWT-токенов. Хранение учётных данных осуществляется в отдельной базе данных \texttt{auth\_db} PostgreSQL, что обеспечивает изоляцию аутентификационных данных от бизнес-данных приложения.

Сервис предоставляет следующие HTTP-эндпоинты~\cite{rfc9110}: \texttt{POST /api/v1/auth/register} для регистрации нового пользователя, \texttt{POST /api/v1/auth/login} для входа и получения токенов, \texttt{POST /api/v1/auth/refresh} для обновления access-токена по refresh-токену, \texttt{GET /api/v1/auth/me} для получения профиля текущего пользователя и \texttt{PUT /api/v1/auth/profile} для обновления имени и аватара.

Пароли пользователей хранятся в хешированном виде. Для хеширования используется алгоритм Argon2id~--- современный победитель конкурса Password Hashing Competition, рекомендованный OWASP в качестве первого выбора для хеширования паролей. Алгоритм устойчив к атакам с применением специализированного оборудования (ASIC, FPGA) благодаря требованиям к памяти.

\subsection{Структура токенов аутентификации}

Для аутентификации API-запросов применяется механизм JSON Web Tokens (JWT)~\cite{rfc7519}. Система выпускает два типа токенов: access-токен с коротким сроком жизни (15 минут) для авторизации запросов, и refresh-токен с длинным сроком жизни (7 дней) для получения нового access-токена без повторного ввода учётных данных.

Полезная нагрузка (claims) JWT-токена содержит следующие поля, представленные структурой \texttt{Claims} в листинге~\ref{lst:jwt-claims}:

\begin{lstlisting}[language=Rust, caption={Структура полезной нагрузки JWT-токена}, label={lst:jwt-claims}]
#[derive(Debug, Serialize, Deserialize, Clone, PartialEq)]
pub struct Claims {
    pub sub: String,         // Subject: user UUID
    pub email: String,       // User email address
    pub role: String,        // Role: user / moderator / admin
    pub exp: i64,            // Expiration timestamp (Unix)
    pub iat: i64,            // Issued-at timestamp (Unix)
    pub token_type: TokenType,
}

#[derive(Debug, Serialize, Deserialize, Clone, PartialEq)]
pub enum TokenType {
    Access,
    Refresh,
}
\end{lstlisting}

Наличие поля \texttt{token\_type} позволяет различать access- и refresh-токены на уровне валидации. Это предотвращает использование refresh-токена в качестве access-токена~--- потенциальную уязвимость, которая присутствует в системах, различающих токены только по сроку жизни.

Метод \texttt{validate\_token} сервиса \texttt{JwtService} выполняет верификацию подписи токена ключом HMAC-SHA256 и проверку срока действия. При истечении срока возвращается специфическая ошибка \texttt{TokenExpired}, которая обрабатывается клиентом для автоматического обновления токена.

\subsection{Промежуточный слой аутентификации}

Защита эндпоинтов сервиса маршрутов реализована через промежуточный слой (middleware) Axum. Middleware извлекает Bearer-токен из заголовка \texttt{Authorization}, валидирует его и внедряет данные аутентифицированного пользователя в расширения запроса (request extensions), делая их доступными для всех обработчиков. Реализация представлена в листинге~\ref{lst:auth-middleware}.

\begin{lstlisting}[language=Rust, caption={Промежуточный слой аутентификации}, label={lst:auth-middleware}]
pub async fn auth_middleware(
    State(state): State<Arc<AppState>>,
    mut request: Request,
    next: Next,
) -> Result<Response, (StatusCode, String)> {
    let token = match request
        .headers()
        .get("Authorization")
        .and_then(|h| h.to_str().ok())
        .and_then(|h| h.strip_prefix("Bearer "))
    {
        Some(t) => t,
        None => return Err((StatusCode::UNAUTHORIZED,
            "Missing Authorization header".to_string())),
    };

    let claims = state.jwt_service.validate_token(token)
        .map_err(|e| (StatusCode::UNAUTHORIZED, e.to_string()))?;

    if claims.token_type != TokenType::Access {
        return Err((StatusCode::UNAUTHORIZED,
            "Invalid token type".to_string()));
    }

    let authenticated_user = AuthenticatedUser {
        user_id: Uuid::parse_str(&claims.sub)?,
        email: claims.email,
        role: claims.role,
    };

    request.extensions_mut().insert(authenticated_user);
    Ok(next.run(request).await)
}
\end{lstlisting}

Структура \texttt{AuthenticatedUser}, внедряемая в расширения запроса, содержит три поля: идентификатор пользователя (\texttt{user\_id} типа UUID), адрес электронной почты и роль. Обработчики защищённых эндпоинтов извлекают эту структуру с помощью экстрактора Axum \texttt{Extension<AuthenticatedUser>}, что гарантирует: обработчик получает управление только при наличии корректного токена.

В текущей реализации используются три системные роли доступа: \texttt{user}, \texttt{moderator} и \texttt{admin}. Роль \texttt{user} является базовой и назначается всем новым аккаунтам; она предоставляет доступ к созданию и использованию собственных маршрутов и социальным функциям. Роль \texttt{moderator} расширяет базовые права возможностью удалять чужие комментарии и маршруты при модерации контента. Роль \texttt{admin} дополнительно открывает административные HTTP-эндпоинты и интерфейс управления пользователями, категориями и статистикой. При этом термины <<гид>> и <<турист>>, используемые далее в работе, обозначают сценарии использования системы, а не отдельные технические роли авторизации.

\section{Реализация сервиса маршрутов}

\subsection{Архитектура и внедрение зависимостей}

Сервис маршрутов является центральным компонентом системы Guide Helper и реализует основную бизнес-логику: управление маршрутами, социальные функции, AI-ассистента, уведомления и администрирование. Сервис разработан на Rust/Axum и следует принципам чистой архитектуры с разделением на слои: delivery (HTTP-обработчики), usecase (бизнес-логика) и repository (работа с базой данных).

Состояние приложения \texttt{AppState} содержит все сервисы бизнес-логики (use cases) и инфраструктурные зависимости. Тип \texttt{AppState} параметризован конкретными реализациями репозиториев через дженерики, что позволяет подменять реализации при тестировании. В листинге~\ref{lst:appstate} представлены ключевые поля структуры.

\begin{lstlisting}[language=Rust, caption={Структура состояния приложения AppState}, label={lst:appstate}]
pub struct AppState {
    pub routes_usecase: RoutesUseCase<PostgresRouteRepository>,
    pub comments_usecase: CommentsUseCase<...>,
    pub chat_usecase: ChatUseCase<...>,
    pub jwt_service: JwtService,
    pub nats_client: Option<async_nats::Client>,
    // WebSocket channels per route_id
    pub ws_channels: Arc<RwLock<HashMap<Uuid, broadcast::Sender<String>>>>,
    pub chat_rate_limits: Arc<RwLock<HashMap<Uuid, (Instant, u32)>>>,
}
\end{lstlisting}

Поле \texttt{ws\_channels} представляет собой словарь каналов рассылки (broadcast channels) библиотеки Tokio, индексированных по идентификатору маршрута. Это позволяет каждому WebSocket-соединению подписаться на уведомления конкретного маршрута без глобальной блокировки.

\subsection{Доменная модель маршрута}

Доменная модель маршрута представлена структурой \texttt{Route} и вложенной структурой \texttt{RoutePoint}. Точки маршрута хранятся в поле типа \texttt{JSONB} в PostgreSQL, что обеспечивает гибкость схемы без дополнительных JOIN-запросов. Каждая точка может содержать необязательные данные фотографии, описываемые структурой \texttt{PhotoData} (листинг~\ref{lst:domain-route}).

\begin{lstlisting}[language=Rust, caption={Доменные модели Route и RoutePoint}, label={lst:domain-route}]
pub struct Route {
    pub id: Uuid,
    pub user_id: Uuid,
    pub name: String,
    pub points: Vec<RoutePoint>,  // stored as JSONB
    pub share_token: Option<Uuid>,
    pub category_ids: Vec<Uuid>,
    pub seasons: Vec<String>,
    pub description: Option<String>,
    pub start_location: Option<String>,
    pub end_location: Option<String>,
    pub created_at: DateTime<Utc>,
    pub updated_at: DateTime<Utc>,
}

pub struct RoutePoint {
    pub lat: f64,
    pub lng: f64,
    pub name: Option<String>,
    pub segment_mode: Option<String>,
    pub photo: Option<PhotoData>,
}

pub struct PhotoData {
    pub original: String,
    pub thumbnail_url: Option<String>,
    pub status: PhotoStatus,  // Pending | Processing | Done | Failed
}
\end{lstlisting}

Для обеспечения обратной совместимости реализован специальный десериализатор поля \texttt{photo}: он принимает как строку в формате data URL (старый формат), так и объект \texttt{PhotoData} (новый формат). Это позволяет поддерживать маршруты, созданные до введения полной обработки фотографий, без миграции данных.

\subsection{Бизнес-логика маршрутов}

Класс \texttt{RoutesUseCase} реализует всю бизнес-логику работы с маршрутами. Базовый конструктор принимает реализацию репозитория \texttt{RouteRepository}; опциональные внешние зависимости~--- клиент Nominatim для геокодирования и vision-клиенты для генерации AI-описаний~--- подключаются через методы-конфигураторы. Такой подход позволяет создавать экземпляры use case без внешних сервисов в юнит-тестах и гибко включать только те интеграции, которые реально нужны в конкретном окружении.

При создании маршрута система выполняет следующую последовательность действий: валидирует входные данные, создаёт объект \texttt{Route} с новым UUID, сохраняет его в PostgreSQL, а затем в фоновой задаче Tokio выполняет обратное геокодирование начальной и конечной точки через Nominatim API. Фоновая задача сохраняет результаты геокодирования в поля \texttt{start\_location} и \texttt{end\_location} маршрута, не блокируя HTTP-ответ пользователю.

Если точки маршрута содержат фотографии в формате base64 data URL, use case создаёт задачу \texttt{PhotoProcessTask} и публикует её в NATS JetStream для асинхронной обработки.

Механизм шеринга маршрутов реализован через уникальный токен типа UUID, хранящийся в поле \texttt{share\_token}. При включении публичного доступа сервис генерирует новый токен; при отключении~--- удаляет его. Публичный маршрут доступен по пути \texttt{/api/v1/shared/\{token\}} без авторизации.

\subsection{Геокодирование и статистика}

Интеграция с Nominatim API обеспечивает автоматическое определение названий начальной и конечной точки маршрута по географическим координатам. Запрос к Nominatim выполняется в фоновой задаче через 100 миллисекунд после создания или обновления маршрута, чтобы не нагружать API при частых обновлениях.

Расчёт статистики маршрута выполняется на стороне сервиса маршрутов с применением алгоритмов, описанных в разделе~1.5. Формула Гаверсинуса применяется для вычисления расстояния между последовательными точками; данные о высотах запрашиваются у Open-Meteo Elevation API и используются для расчёта суммарного набора высоты по формуле Нейсмита. На основе совокупного расстояния и набора высоты алгоритм классификации автоматически присваивает маршруту уровень сложности: «Лёгкий», «Средний» или «Сложный».

\section{Асинхронная обработка фотографий}

\subsection{Конвейер обработки через NATS JetStream}

Асинхронная обработка фотографий реализована как отдельный микросервис \texttt{photo-worker}, взаимодействующий с основным сервисом маршрутов через брокер сообщений NATS JetStream. Такая архитектура обеспечивает изоляцию отказов: сбой в обработке изображений не влияет на работу API.

При загрузке фотографии пользователь передаёт её в формате base64 data URL вместе с данными точки маршрута. Сервис маршрутов сохраняет точку в базе данных со статусом \texttt{Pending} и публикует задачу в поток NATS JetStream (листинг~\ref{lst:nats-publish}).

\begin{lstlisting}[language=Rust, caption={Публикация задачи обработки фото в NATS}, label={lst:nats-publish}]
pub struct PhotoProcessTask {
    pub route_id: Uuid,
    pub user_id: Uuid,
    pub point_indices: Vec<usize>,
}

impl PhotoProcessTask {
    pub fn from_route(route: &Route) -> Option<Self> {
        let indices: Vec<usize> = route.points
            .iter()
            .enumerate()
            .filter_map(|(i, point)| {
                point.photo.as_ref()
                    .filter(|p| p.original.starts_with("data:"))
                    .map(|_| i)
            })
            .collect();

        if indices.is_empty() { return None; }
        Some(Self { route_id: route.id,
                    user_id: route.user_id,
                    point_indices: indices })
    }
}
\end{lstlisting}

Метод \texttt{from\_route} определяет, какие точки маршрута содержат необработанные фотографии (в формате data URL), и включает их индексы в задачу. Уже обработанные фотографии (с URL вида \texttt{/photos/...}) в задачу не включаются, что обеспечивает идемпотентность при повторной публикации.

\subsection{Обработчик фотографий photo-worker}

Сервис \texttt{photo-worker} при запуске подключается к PostgreSQL, MinIO~\cite{minio} (S3-совместимое объектное хранилище) и NATS JetStream, после чего создаёт постоянного потребителя (durable consumer) потока \texttt{PHOTOS} с политикой \texttt{WorkQueue}. Политика WorkQueue гарантирует, что каждое сообщение будет доставлено ровно одному потребителю и удалено из потока после подтверждения (ACK).

Конвейер обработки одной фотографии включает следующие этапы:
\begin{enumerate}
  \item декодирование base64 data URL в массив байтов;
  \item загрузка изображения в память и масштабирование до максимальной ширины 1920 пикселей с использованием фильтра Ланцоша (Lanczos3) для сохранения качества при уменьшении;
  \item генерация миниатюры с шириной 300 пикселей для быстрой загрузки маркеров на карте;
  \item кодирование в формат JPEG с качеством 85\%~--- оптимальный баланс размера файла и визуального качества;
  \item загрузка оригинала и миниатюры в MinIO по ключам \texttt{\{user\_id\}/\{route\_id\}/photo\_\{index\}.jpg} и \texttt{.../thumb\_\{index\}.jpg};
  \item обновление записи в PostgreSQL: замена base64 на URL файлов MinIO, установка статуса \texttt{Done};
  \item публикация события завершения в NATS по теме \texttt{photos.completed.\{route\_id\}} для уведомления фронтенда через WebSocket.
\end{enumerate}

Реализация функции декодирования и сжатия приведена в листинге~\ref{lst:photo-processing}.

\begin{lstlisting}[language=Rust, caption={Декодирование data URL и сжатие изображения}, label={lst:photo-processing}]
pub fn decode_data_url(data_url: &str) -> Result<Vec<u8>> {
    let comma_pos = data_url.find(',')
        .ok_or_else(|| anyhow!("invalid data URL"))?;
    let encoded = &data_url[comma_pos + 1..];
    base64::engine::general_purpose::STANDARD
        .decode(encoded)
        .context("failed to decode base64")
}

pub fn compress_image(data: &[u8], max_width: u32) -> Result<Vec<u8>> {
    let img = image::load_from_memory(data)?;
    let img = if img.width() > max_width {
        let ratio = max_width as f64 / img.width() as f64;
        let new_h = (img.height() as f64 * ratio) as u32;
        img.resize(max_width, new_h, FilterType::Lanczos3)
    } else { img };

    let mut buf = Cursor::new(Vec::new());
    img.write_to(&mut buf, ImageFormat::Jpeg)?;
    Ok(buf.into_inner())
}
\end{lstlisting}

\subsection{Уведомления в реальном времени}

После завершения обработки фотографии \texttt{photo-worker} публикует событие в NATS по теме \texttt{photos.completed.\{route\_id\}}. Сервис маршрутов подписан на эти события и при получении рассылает уведомление по WebSocket-каналу соответствующего маршрута.

WebSocket-соединение~\cite{rfc6455} устанавливается клиентом при открытии редактора маршрута через эндпоинт \texttt{GET /api/v1/routes/\{id\}/ws}. Для каждого маршрута создаётся канал рассылки Tokio (\texttt{broadcast::Sender}), хранящийся в поле \texttt{ws\_channels} объекта \texttt{AppState}. Все клиенты, просматривающие один и тот же маршрут, получают уведомление одновременно.

Уведомление содержит идентификатор точки, URL обработанного оригинала и URL миниатюры. Клиент обновляет маркер на карте: заменяет индикатор загрузки на миниатюру фотографии без перезагрузки страницы.

\section{Интеграция ИИ-ассистента}

\subsection{Многопровайдерная архитектура}

Система Guide Helper поддерживает несколько провайдеров языковых моделей: OpenAI-совместимые API, внешний шлюз \texttt{claude-code-api} для работы с Claude в OpenAI-совместимом режиме, прямой Anthropic Claude API, а также локальный Ollama. Выбор активного провайдера осуществляется через переменную окружения \texttt{AI\_PROVIDER} или динамически через систему флагов функциональности Unleash.

Конфигурация провайдеров задаётся в структуре \texttt{AppConfig}:
\begin{itemize}
  \item \texttt{openai\_api\_key}, \texttt{openai\_base\_url}, \texttt{openai\_model}~--- настройки OpenAI-совместимого API;
  \item \texttt{anthropic\_api\_key}, \texttt{anthropic\_model}~--- настройки Anthropic;
  \item \texttt{claude\_base\_url}, \texttt{claude\_api\_key}, \texttt{claude\_model}~--- настройки внешнего \texttt{claude-code-api} или иного OpenAI-совместимого шлюза к Claude;
  \item \texttt{ollama\_chat\_base\_url}, \texttt{ollama\_chat\_model}, \texttt{ollama\_base\_url}, \texttt{ollama\_vision\_model}~--- настройки локального Ollama;
  \item \texttt{unleash\_url}, \texttt{unleash\_api\_token}~--- опциональная интеграция с Unleash для A/B-переключения провайдеров.
\end{itemize}

Все провайдеры реализуют единый трейт \texttt{AiChatClient}, что позволяет подменять их без изменения кода бизнес-логики. Класс \texttt{DynamicAiClient} реализует тот же трейт и делегирует вызовы активному провайдеру, определяемому в runtime на основе флага Unleash или статической конфигурации.

\subsection{Диалоговый интерфейс с вызовом инструментов}

ИИ-ассистент реализован как полноценный диалоговый интерфейс с поддержкой вызова инструментов (function calling / tool use). При получении сообщения от пользователя сервис Chat отправляет историю диалога и доступные инструменты провайдеру ЯМ. Если модель решает вызвать инструмент, сервис выполняет соответствующую операцию и возвращает результат обратно модели для формирования финального ответа.

Реализованы следующие инструменты ассистента:
\begin{itemize}
  \item geocode~--- поиск географических координат по текстовому запросу через Nominatim API;
  \item search\_routes~--- поиск маршрутов в каталоге по тексту, категории и местоположению;
  \item get\_route\_details~--- получение подробной информации о конкретном маршруте по UUID;
  \item get\_weather~--- запрос прогноза погоды для заданных координат через Open-Meteo~\cite{openmeteo}.
\end{itemize}

Цикл вызова инструментов ограничен параметром \texttt{chat\_max\_tool\_iterations} (по умолчанию 5 итераций), что предотвращает бесконечные циклы. Для предотвращения злоупотреблений применяется оконное ограничение частоты запросов: не более \texttt{chat\_rate\_limit\_max} (10) сообщений за \texttt{chat\_rate\_limit\_window\_secs} (60 секунд) для одного пользователя.

Генерация текстовых описаний маршрута по фотографиям выполняется через vision-провайдер, настроенный для сервиса маршрутов. В текущей реализации поддерживаются Anthropic Claude и локальный Ollama; при наличии ключа Anthropic приоритетно используется Claude, в противном случае~--- Ollama. Сервис отправляет модели URL обработанных фотографий маршрута и запрос на генерацию описания на русском языке объёмом 3--5 предложений.

\section{Клиентское веб-приложение}

\subsection{Структура клиентского веб-приложения}

Клиентская часть системы разработана с использованием библиотеки React~19 и языка TypeScript. Приложение организовано по страничной архитектуре: каждая страница соответствует отдельному разделу функциональности системы.

Основные страницы приложения:
\begin{itemize}
  \item MapPage~--- главная страница с интерактивной картой и редактором маршрутов;
  \item ExplorePage~--- каталог публичных маршрутов с поиском и фильтрацией;
  \item ProfilePage~--- профиль пользователя со списком созданных маршрутов;
  \item BookmarksPage~--- список маршрутов, добавленных в избранное;
  \item ChatPage~--- интерфейс ИИ-ассистента;
  \item AdminPage~--- административная панель управления;
  \item SharedRoutePage~--- публичный просмотр маршрута по уникальной ссылке.
\end{itemize}

Приложение поддерживает светлую и тёмную темы оформления через переключатель в навигационной панели. Интернационализация реализована с помощью библиотеки i18next: поддерживаются русский и английский языки, выбор языка сохраняется в localStorage браузера.

\subsection{Интерактивная карта на Leaflet}

Интерактивная карта реализована на основе библиотеки Leaflet~\cite{leaflet} с обёрткой React-Leaflet~\cite{reactleaflet}. Система поддерживает несколько источников тайлов: OpenStreetMap~\cite{openstreetmap} (основной), Yandex Maps и исторический слой OpenHistoricalMap. Тайлы кешируются сервисом \texttt{cache} через Redis~\cite{redis} для снижения нагрузки на внешние API и ускорения загрузки карты.

Компонент редактора маршрутов (\texttt{RouteBuilder}) поддерживает два режима добавления точек: ручной (клик по карте) и автоматический (маршрутизация через OSRM~\cite{osrm} с построением кратчайшего пути по дорогам). Между последовательными точками отображается линия маршрута; для каждого сегмента можно выбрать режим передвижения (пешком, велосипед, автомобиль).

Точки с прикреплёнными фотографиями отображаются на карте специальными маркерами с миниатюрами. При клике на маркер открывается всплывающее окно с полноразмерной фотографией, названием и описанием точки. Компонент профиля высот (\texttt{ElevationProfile}) отображает интерактивный график изменения высоты вдоль маршрута; при наведении курсора на точку графика соответствующая точка подсвечивается на карте.

\subsection{Прогрессивное веб-приложение и настольная версия}

Клиентское приложение сконфигурировано как прогрессивное веб-приложение: оно регистрирует Service Worker, который кеширует ресурсы приложения и ранее загруженные маршруты для работы в офлайн-режиме. Пользователь может добавить приложение на главный экран мобильного устройства.

Для создания настольной версии (Windows, macOS, Linux) используется фреймворк Tauri~\cite{tauri}. Tauri оборачивает веб-приложение в нативную оболочку с доступом к файловой системе операционной системы, при этом размер итогового исполняемого файла значительно меньше, чем у Electron-приложений, благодаря использованию системного WebView и ядра на Rust.

Компоненты \texttt{WeatherWidget} и \texttt{NotificationPanel} обращаются к соответствующим API-эндпоинтам для отображения прогноза погоды для района маршрута и уведомлений о новых комментариях, отметках «нравится» и оценках в реальном времени.

\section{Схема базы данных и репозиторный слой}

Сервис маршрутов использует отдельную базу данных \texttt{routes\_db} PostgreSQL, содержащую 11 таблиц. Доступ к данным реализован через репозиторный слой~\cite{fowler2018}: каждый агрегат (маршруты, комментарии, лайки и т.д.) имеет собственный трейт-интерфейс и конкретную реализацию \texttt{PostgresXxxRepository}. Такое разделение позволяет подменять реализацию в тестах.

Ключевые таблицы базы данных:
\begin{itemize}
  \item \texttt{routes}~--- маршруты пользователей; поле \texttt{points} типа \texttt{JSONB} хранит массив точек с координатами и данными фотографий; поле \texttt{share\_token} содержит опциональный UUID~\cite{rfc4122} публичной ссылки;
  \item \texttt{route\_categories}~--- таблица связи «многие-ко-многим» между маршрутами и категориями;
  \item \texttt{comments}, \texttt{route\_likes}, \texttt{route\_ratings}~--- социальные функции; в таблице рейтингов пара (\texttt{route\_id}, \texttt{user\_id}) уникальна, что гарантирует не более одной оценки от одного пользователя;
  \item \texttt{notifications}~--- уведомления о событиях (комментарий, лайк, оценка) с полем \texttt{is\_read};
  \item \texttt{chat\_messages}~--- история диалогов с ИИ-ассистентом; сообщения объединяются по \texttt{conversation\_id}, а роль сообщения (\texttt{user} / \texttt{assistant}) и опциональный массив вызовов инструментов хранятся в \texttt{JSONB};
  \item \texttt{route\_bookmarks}~--- избранные маршруты; пара (\texttt{user\_id}, \texttt{route\_id}) уникальна.
\end{itemize}

Для получения данных применяется библиотека sqlx~\cite{sqlx} с проверкой SQL-запросов на этапе компиляции: \texttt{sqlx::query\_as!} сопоставляет результат запроса с Rust-структурой при сборке, что исключает несоответствие типов в runtime. Поддержка типа \texttt{JSONB} реализована через атрибут \texttt{\#[sqlx(json)]}, автоматически выполняющий сериализацию и десериализацию через \texttt{serde\_json}.

\section{Инфраструктура автоматической сборки, проверки и развёртывания}

\subsection{Контейнеризация через Docker Compose}

Для локальной разработки и тестирования все компоненты системы оркестрированы через Docker Compose~\cite{docker}. Файл \texttt{docker-compose.yml} описывает десять сервисов: PostgreSQL (с инициализацией двух баз данных \texttt{auth\_db} и \texttt{routes\_db}), NATS JetStream, MinIO, Redis, сервис аутентификации, сервис маршрутов, обработчик фотографий, сервис кеширования тайлов, сервис проксирования тайлов и фронтенд.

Зависимости между сервисами обеспечиваются директивой \texttt{depends\_on} с проверками готовности (\texttt{healthcheck}): сервис маршрутов ожидает готовности PostgreSQL (через \texttt{pg\_isready}), NATS (через HTTP-мониторинг на порту 8222) и сервиса аутентификации. Таблица~\ref{tab:docker-services} содержит перечень контейнеров с назначенными портами.

\begin{table}[H]
\caption{Сервисы Docker Compose}
\label{tab:docker-services}
\begin{tabularx}{\textwidth}{|l|X|c|}
\hline
\textbf{Сервис} & \textbf{Назначение} & \textbf{Порт} \\
\hline
postgres & PostgreSQL 16, базы auth\_db и routes\_db & 5433 \\
\hline
nats & NATS JetStream, брокер сообщений & 4222 \\
\hline
minio & MinIO S3, хранилище фотографий & 9000 \\
\hline
redis & Redis 7, кеш тайлов карты & 6379 \\
\hline
auth & Сервис аутентификации (Rust/Axum) & 8086 \\
\hline
routes & Сервис маршрутов (Rust/Axum) & 8088 \\
\hline
photo-worker & Обработчик фотографий (Rust/NATS) & --- \\
\hline
cache & Кеширование тайлов (Go/Gin) & 8085 \\
\hline
tiles & Проксирование тайлов (Go/Gin) & 8087 \\
\hline
frontend & React-приложение & 3005 \\
\hline
\end{tabularx}
\end{table}

\subsection{Kubernetes-развёртывание}

Для продуктивного развёртывания системы подготовлены манифесты Kubernetes~\cite{luksa2018,kubernetes}, управляемые через Kustomize. Каждый микросервис оформлен как отдельный \texttt{Deployment} с соответствующим \texttt{Service}. Конфиденциальные данные (JWT-секрет, ключи MinIO, токены API) хранятся в объектах \texttt{Secret}.

Для автоматического горизонтального масштабирования применяется \texttt{HorizontalPodAutoscaler}: при превышении порога загрузки CPU Kubernetes автоматически увеличивает число реплик сервиса маршрутов. Обработчик фотографий \texttt{photo-worker} допускает независимое масштабирование для ускорения обработки при высокой нагрузке.

\subsection{Непрерывная интеграция и доставка}

Конвейер автоматической сборки, проверки и развёртывания реализован средствами GitHub Actions~\cite{githubactions}. При каждом коммите в ветку \texttt{main} запускается рабочий процесс, включающий следующие этапы:
\begin{enumerate}
  \item запуск юнит-тестов для сервиса маршрутов (\texttt{cargo test});
  \item сборка Docker-образов для всех микросервисов и фронтенда;
  \item публикация образов в реестр GitHub Container Registry (\texttt{ghcr.io}) с тегом, соответствующим хешу коммита;
  \item обновление тегов образов в Kubernetes-манифестах и применение изменений в кластере.
\end{enumerate}

Образы собираются с применением многоэтапной сборки (\texttt{multi-stage build}): финальный образ Rust-сервисов основан на \texttt{rust:1.91.0-alpine} и содержит только скомпилированный бинарный файл без исходного кода, что минимизирует размер и поверхность атаки.

\section{Руководство пользователя и администратора}

\subsection{Системные требования и развёртывание}

Для развёртывания системы Guide Helper необходимо следующее аппаратное и программное обеспечение:
\begin{itemize}
  \item операционная система: Linux (Ubuntu~22.04/Debian~12), macOS~13+, Windows~10+ с WSL~2;
  \item Docker Engine 24.0 или выше, Docker Compose 2.20 или выше~\cite{docker};
  \item оперативная память: не менее 4~ГБ (рекомендуется 8~ГБ);
  \item дисковое пространство: не менее 10~ГБ (для образов и данных);
  \item сетевое подключение к интернету (для загрузки образов и картографических тайлов).
\end{itemize}

Порядок установки:
\begin{enumerate}
  \item Клонировать репозиторий:
\begin{lstlisting}[language=bash, caption={Клонирование репозитория}]
git clone https://github.com/jaennil/guide_helper.git
cd guide_helper
\end{lstlisting}
  \item Создать файл конфигурации из шаблона и заполнить переменные окружения:
\begin{lstlisting}[language=bash, caption={Настройка конфигурации}]
cp .env.example .env
# Обязательные параметры:
# JWT_SECRET    -- секретный ключ для подписи JWT (минимум 32 символа)
# MINIO_ROOT_PASSWORD -- пароль MinIO (минимум 8 символов)
\end{lstlisting}
  \item Запустить все сервисы командой:
\begin{lstlisting}[language=bash, caption={Запуск системы}]
docker compose up -d
\end{lstlisting}
  \item После запуска веб-приложение доступно по адресу \texttt{http://localhost:5173}.
\end{enumerate}

\subsection{Руководство пользователя (роль «Гид»)}

Гид~--- основной пользователь системы, создающий туристические маршруты. Порядок работы:

\textbf{Регистрация и вход.} Открыть страницу входа по адресу приложения. Нажать кнопку <<Регистрация>>, заполнить поля <<Имя>>, <<Email>> и <<Пароль>> (минимум 8~символов). После успешной регистрации выполнить вход.

\textbf{Создание маршрута.}
\begin{enumerate}
  \item На главной странице нажать кнопку <<Новый маршрут>>.
  \item Присвоить маршруту название, выбрать категорию из предложенного списка.
  \item На интерактивной карте добавить точки маршрута: кликнуть на карту~--- точка добавляется автоматически.
  \item Для каждой точки можно загрузить фотографию: нажать на маркер, выбрать файл с диска (JPEG/PNG, до 10~МБ). Система автоматически считывает GPS-координаты из EXIF-метаданных и привязывает фотографию к точке. При отсутствии EXIF используются координаты клика.
  \item Нажать кнопку <<Построить маршрут>>~--- система автоматически прокладывает оптимальный маршрут между точками через сервис OSRM с учётом дорог и троп~\cite{osrm}.
  \item В правой панели отображается статистика: расстояние, набор и сброс высоты, расчётное время прохождения по формуле Нейсмита.
  \item Сохранить маршрут кнопкой <<Сохранить>>.
\end{enumerate}

\textbf{Публикация и шеринг.}
\begin{enumerate}
  \item Открыть сохранённый маршрут из раздела <<Мои маршруты>> в профиле.
  \item Нажать кнопку <<Опубликовать>>~--- маршрут появится в публичном каталоге.
  \item Для передачи маршрута клиентам нажать <<Поделиться>>: система формирует уникальную ссылку. По этой ссылке любой пользователь (без регистрации) может просмотреть интерактивную карту маршрута с фотографиями.
  \item Маршрут также доступен для экспорта в форматах GPX (для GPS-навигаторов) и KML (для Google Earth).
\end{enumerate}

\textbf{ИИ-ассистент.} На любой странице приложения доступна кнопка вызова ИИ-ассистента (значок чата в правом нижнем углу). Ассистент может: геокодировать адреса и находить места на карте, подбирать маршруты по критериям (сложность, расстояние, категория), предоставлять прогноз погоды для точек маршрута.

\subsection{Руководство пользователя (роль «Турист»)}

Турист~--- пользователь, просматривающий и оценивающий маршруты. Функции доступны без регистрации (просмотр) и с регистрацией (социальные функции).

\textbf{Просмотр каталога.} Открыть раздел <<Каталог>> (значок компаса в навигации). На странице отображается список опубликованных маршрутов с обложками, рейтингом и статистикой. Доступна фильтрация по категориям, сложности и сезону; сортировка по дате, рейтингу и количеству отметок «нравится»; поиск по названию и описанию.

\textbf{Просмотр маршрута.} Нажать на карточку маршрута. Откроется интерактивная карта с маршрутом и маркерами фотографий. Клик по маркеру~--- всплывающее окно с фотографией и описанием точки. Под картой отображается профиль высот и полная статистика маршрута.

\textbf{Социальные функции (требуется авторизация).}
\begin{itemize}
  \item Отметить маршрут кнопкой <<Нравится>> (сердечко).
  \item Добавить в закладки кнопкой <<Сохранить>> для быстрого доступа из профиля.
  \item Оставить рейтинг (1--5~звёзд) и текстовый комментарий.
\end{itemize}

\subsection{Руководство администратора}

Администратор получает расширенный доступ через ту же форму входа. Роль назначается в базе данных (поле \texttt{role = 'admin'}).

\textbf{Панель администратора} доступна по ссылке <<Администрирование>> в меню профиля. Функции:

\begin{itemize}
  \item Управление пользователями: просмотр списка пользователей и изменение роли (\texttt{user}~/ \texttt{moderator}~/ \texttt{admin}).
  \item Модерация контента: просмотр маршрутов и комментариев, удаление нарушающего контента.
  \item Управление категориями: создание, редактирование и удаление категорий маршрутов.
  \item Статистика системы: общее число пользователей, маршрутов и комментариев.
\end{itemize}

\textbf{Роль модератора.} В пользовательском интерфейсе модератор не имеет отдельной административной панели, однако на уровне серверной авторизации получает расширенные права на удаление чужих маршрутов и комментариев при модерации контента.

\textbf{Мониторинг.} Логи всех микросервисов доступны командой:
\begin{lstlisting}[language=bash, caption={Просмотр логов системы}]
docker compose logs -f --tail=100
\end{lstlisting}

\section*{Выводы по второй главе}
\addcontentsline{toc}{section}{Выводы по второй главе}

В ходе второй главы реализованы все ключевые компоненты системы Guide Helper.
Сервис аутентификации обеспечивает выдачу и обновление JWT access/refresh токенов
с хранением паролей по алгоритму Argon2id. Сервис маршрутов реализует полный
Полный цикл операций создания, чтения, изменения и удаления маршрутов с точками, поддерживает импорт GeoJSON, экспорт GPX/KML
и расчёт статистик на основе формулы Хаверсина и правила Найсмита.

Конвейер асинхронной обработки фотографий реализован через NATS JetStream: задачи
обработки сохраняются с гарантией доставки, фотографии сжимаются и масштабируются
с помощью фильтра Ланцоша, результаты сохраняются в MinIO, а клиент получает
уведомление через WebSocket. Полный цикл обработки изображения до 5~МБ занимает
2--4~секунды, что соответствует нефункциональным требованиям.

ИИ-ассистент реализован с поддержкой нескольких провайдеров языковых моделей и
вызовом инструментов (геокодирование, поиск маршрутов, прогноз погоды). Клиентское
React-приложение включает семь страниц с интерактивной картой на Leaflet,
поддержкой тёмной темы и интернационализацией. Настроен конвейер автоматической сборки, проверки и развёртывания на
GitHub Actions с автоматическим тестированием, сборкой Docker-образов и
развёртыванием в Kubernetes-кластер.
